\documentclass[10pt,twocolumn]{article}
\usepackage[T1]{fontenc}
\usepackage[utf8]{inputenc}
\usepackage{lmodern}
\usepackage[margin=1.8cm,top=2.4cm,bottom=2cm,columnsep=0.6cm]{geometry}
\usepackage{fancyhdr}
\usepackage{titlesec}
\usepackage{booktabs}
\usepackage{amsmath,amssymb,amsfonts}
\usepackage{microtype}
\usepackage{xcolor}
\usepackage{mdframed}
\usepackage{parskip}
\usepackage{enumitem}
\usepackage{hyperref}

\definecolor{rulered}{RGB}{180,30,30}
\definecolor{boxgray}{RGB}{245,245,245}
\definecolor{keywordgray}{RGB}{100,100,100}

\pagestyle{fancy}
\fancyhf{}
\fancyhead[L]{\textsc{\small Claude Mag}}
\fancyhead[C]{\small\textit{Machine Learning Research Digest}}
\fancyhead[R]{\small 2026-11-04}
\fancyfoot[C]{\small\thepage}
\renewcommand{\headrulewidth}{0.8pt}

\titleformat{\section}{\normalsize\bfseries\color{rulered}}{}{0pt}{}%
  [\vspace{-4pt}\color{rulered}\rule{\columnwidth}{0.4pt}]
\titlespacing{\section}{0pt}{8pt}{4pt}

\hypersetup{colorlinks=true,urlcolor=rulered,linkcolor=black}

\begin{document}
\setlength{\parindent}{0pt}
\setlength{\parskip}{4pt}

{\small\color{keywordgray}\textbf{Keywords:}
  Transformer theory \textbullet{}
  attention low-rank structure \textbullet{}
  random matrix theory \textbullet{}
  softmax spectral analysis}\par\vspace{4pt}

{\LARGE\bfseries\raggedright
  Spectral Degeneration of Self-Attention}\par\vspace{4pt}

{\large\itshape\raggedright
  Under Isotropic Score Geometry, Softmax Collapses the Attention
  Spectrum; Causal Masking Grows Stable Rank as $O(\log L)$}\par\vspace{6pt}

{\small\textbf{By} Mengda Li, Jianfeng Yao}\par\vspace{2pt}

{\small\color{keywordgray}arXiv:2611.02778 \textbullet{}
  cs.LG / stat.ML \textbullet{} NeurIPS 2026 \textbullet{}
  CUHK Shenzhen}
\par\vspace{2pt}\noindent{\color{rulered}\rule{\columnwidth}{0.6pt}}\vspace{6pt}

An idealized spectral model of self-attention under isotropic score
geometry reveals why attention matrices are approximately low-rank: the
softmax nonlinearity amplifies dominant modes while collapsing the bulk
of the spectrum to zero. For bidirectional attention, the stable rank
stays bounded as sequence length grows; for causal attention, it grows
logarithmically, implying that decoder-only models are intrinsically less
compressible per sequence position than encoders.

\begin{mdframed}[backgroundcolor=boxgray,linecolor=rulered,linewidth=1pt,
    innerleftmargin=6pt,innerrightmargin=6pt,
    innertopmargin=6pt,innerbottommargin=6pt]
\textbf{\small AT A GLANCE}\par\vspace{4pt}
\begin{description}[leftmargin=0pt,labelsep=4pt,itemsep=2pt,parsep=0pt,
    font=\small\bfseries]
  \item[Problem:] Attention matrices are empirically low-rank, but the
    theoretical reason for this degeneration is unknown.
  \item[Why it matters:] Understanding why attention collapses informs
    principled low-rank compression, regularization design, and the
    choice between encoder and decoder architectures.
  \item[Method:] Idealized Gaussian random matrix model for the score
    matrix, combined with free probability theory to track how softmax
    distorts the eigenspectrum.
  \item[Result:] Softmax amplifies dominant eigenvalues while collapsing
    the bulk; stable rank is $O(1)$ for bidirectional and $O(\log L)$ for
    causal attention.
  \item[Evidence:] Analytical derivations validated on pretrained GPT-2
    and BERT attention heads and synthetic random attention matrices.
\end{description}
\end{mdframed}

\section{The Big Picture}

Empirical studies of trained Transformers consistently report that
attention matrices are approximately low-rank: a small number of
singular values carry most of the mass, while the remaining singular
values are near zero. Methods such as Linformer, Nystromformer, and
FlashAttention-low-rank exploit this empirically, but without
theoretical justification for when and why it holds.

Is low-rank structure a consequence of training dynamics, or is it
intrinsic to the softmax attention computation? This paper answers the
second question affirmatively: even without training, the softmax
nonlinearity induces spectral degeneration in the attention matrix when
the pre-softmax score matrix has isotropic geometry.

\section{Why Existing Theory Falls Short}

\textbf{Empirical studies} document low-rank structure but do not
explain its origin; they attribute it to gradient dynamics without
analysing the effect of the softmax function itself.

\textbf{Linear attention theory} replaces softmax with a kernel
approximation and analyzes the resulting linear map, but this
fundamentally changes the operation and does not explain the behavior
of softmax attention.

\textbf{Rank-collapse results} from overparameterized network theory
apply to weight matrices in MLP layers and do not transfer to the
compositional query--key product structure of attention.

\section{The Isotropic Score Model}

Let $S \in \mathbb{R}^{L \times L}$ be the pre-softmax score matrix
with $S_{ij} = q_i^\top k_j / \sqrt{d}$. The isotropic assumption
models $S$ as a Wigner-type random matrix:
\[
  S = \frac{1}{\sqrt{L}} W + \lambda \mathbf{u}\mathbf{v}^\top
\]
where $W$ has i.i.d.\ entries with mean 0 and variance 1, $\lambda$ is
the signal strength, and $\mathbf{u},\mathbf{v}$ represent a planted
dominant mode. This models a ``generic'' attention head in which most
query--key pairs are unrelated ($W$ term) but a few share a strong
structural alignment ($\lambda$ term).

\section{Main Theoretical Results}

\textbf{Softmax Spectral Transformation.} Let
$A = \mathrm{softmax}(S)$. Under the isotropic model, in the limit
$L \to \infty$ with fixed $d$, the empirical spectral distribution of
$A$ converges to a law concentrated near zero, with an isolated outlier
eigenvalue at:
\[
  \lambda_1(A) \approx \frac{e^{\lambda/\sqrt{L}}}
    {\mathbf{1}^\top e^{S\mathbf{v}} / L}
\]
while the bulk of the spectrum satisfies $\lambda_i(A) = O(1/L)$ for
$i \geq 2$. The \textbf{stable rank} $\rho(A) = \|A\|_F^2 / \|A\|_2^2$
therefore obeys:
\[
  \rho(A) = O(1) \quad \text{(bidirectional attention)}
\]

\textbf{Causal Attention.} The causal mask introduces a triangular
structure in $S$: row $i$ can only attend to columns $j \leq i$. The
row-normalization of softmax then has a gradient: earlier rows normalize
over fewer tokens. The analysis shows that this asymmetry causes the
stable rank to grow as:
\[
  \rho(A_{\text{causal}}) = O(\log L)
\]

\section{Proof Overview}

The derivation proceeds in three steps:
\begin{enumerate}[itemsep=2pt,parsep=0pt]
  \item Apply the semicircle law from random matrix theory to
    characterize the eigenspectrum of $S$.
  \item Use free probability calculus to track the distortion of the
    spectral measure through the entrywise exponential
    $\exp(S)$.
  \item Show that the row-normalization step of softmax is approximated
    by a rank-1 update to $\exp(S)$ in the high-dimensional limit, and
    derive its effect on singular values.
\end{enumerate}

For causal attention, the proof additionally shows that the
per-row normalization constants form a monotonically increasing sequence
in the row index, leading to a non-uniform reweighting of singular
vectors that grows the stable rank logarithmically.

\section{What the Guarantee Says}

The spectral collapse results are asymptotic ($L \to \infty$, fixed $d$)
and require the isotropic geometry assumption, which may not hold for
specialized attention heads in trained models. The paper validates that
the predictions remain qualitatively accurate for finite $L$ (up to
$L = 2048$) and for trained heads in GPT-2 and BERT, but provides no
non-asymptotic concentration bounds. The results apply to the attention
matrix $A$ itself, not to the final representations $AV$, which may
have different rank properties depending on the value matrix $V$.

\section{Empirical Validation}

\begin{itemize}[itemsep=2pt,parsep=0pt]
  \item \textbf{Synthetic:} Stable rank under the isotropic model
    matches analytical predictions at $L \in \{128, 512, 2048\}$.
  \item \textbf{GPT-2 (causal):} Empirical stable rank grows as
    $1.3 \log L + 0.7$ across heads, consistent with the $O(\log L)$
    prediction.
  \item \textbf{BERT (bidirectional):} Empirical stable rank stays in
    the range $[1.8, 4.2]$ across heads and all $L$ tested, consistent
    with the $O(1)$ bound.
  \item \textbf{Head specialization:} The planted outlier $\lambda$
    correctly predicts the magnitude of the dominant singular value
    across heads with varying degrees of specialization.
\end{itemize}

\section{Implications for Practice}

\textbf{Low-rank compression:} The $O(1)$ stable rank for bidirectional
attention justifies using rank-$r$ approximations with $r$ independent
of $L$. For causal models, rank should scale mildly with $L$.

\textbf{Architecture choice:} Encoder models can achieve the same
compression ratio regardless of context length. Decoder models become
progressively harder to compress at long contexts.

\textbf{Regularization:} Encouraging isotropic score geometry during
training via spectral normalization of $Q$ and $K$ projections may
stabilize the spectral structure and improve compressibility.

\section*{Reference}

Mengda Li, Jianfeng Yao.
\textbf{Spectral Degeneration of Self-Attention under Isotropic Score
Geometry.}
arXiv:2611.02778. NeurIPS 2026.
\url{https://arxiv.org/abs/2611.02778}

\end{document}
